\chapter{Thermodynamics and Equation of State}

\section{First Law and Heat Capacities}

\subsection{Fundamental Identity}
\begin{theory}
For systems with variable particle number,
\begin{equation}
    dQ = dU + dW - \mu dN,
\end{equation}
with $dW=P\,dV$. For closed systems ($dN=0$),
\begin{equation}
    dQ = dU + P\,dV.
\end{equation}
\end{theory}

\begin{importantnote}
Sign conventions matter. In these notes, $dQ>0$ means heat into the system and $dW=P\,dV$ is work done by the system.
\end{importantnote}

\subsection{Adiabatic, Isochoric, Isobaric, Isothermal Processes}
\begin{theory}
\begin{itemize}
    \item Adiabatic ($dQ=0$):
    \begin{equation}
        dU=-P\,dV
    \end{equation}
    \item Isochoric ($dV=0$):
    \begin{equation}
        C_V = \left(\frac{\partial U}{\partial T}\right)_V
    \end{equation}
    \item Isobaric ($P=\text{const}$, ideal gas):
    \begin{equation}
        C_P = C_V + nR
    \end{equation}
    \item Isothermal (ideal gas):
    \begin{equation}
        dU=0,\quad dQ=dW=\frac{nRT}{V}dV
    \end{equation}
\end{itemize}
\end{theory}

\section{Entropy and Thermodynamic Potentials}

\subsection{Entropy and Fundamental Relation}
\begin{theory}
For reversible processes,
\begin{equation}
    dS = \frac{dQ_{\text{rev}}}{T},
\end{equation}
leading to
\begin{equation}
    dU = T\,dS - P\,dV + \mu\,dN.
\end{equation}
\end{theory}

\subsection{Helmholtz Free Energy and Maxwell Relations}
\begin{foundation}
Define $F=U-TS$. Then
\begin{equation}
    dF=-S\,dT-P\,dV+\mu\,dN,
\end{equation}
which gives, for example,
\begin{equation}
    \left(\frac{\partial S}{\partial V}\right)_{T,N}=\left(\frac{\partial P}{\partial T}\right)_{V,N}.
\end{equation}
\end{foundation}

\begin{remark}
For ideal gas, $(\partial U/\partial V)_T=0$, so internal energy depends only on $T$.
This fails for interacting gases such as Van der Waals fluids.
\end{remark}

\section{Photon Gas and Blackbody Radiation}

\subsection{Photon Statistics and EOS}
\begin{theory}
Photons are bosons with $\mu=0$ and $\varepsilon=pc$.
With Bose--Einstein statistics,
\begin{equation}
    \langle n_p\rangle = \frac{1}{e^{pc/k_B T}-1}.
\end{equation}
The energy density is
\begin{equation}
    u=aT^4,\quad a=\frac{8\pi^5k_B^4}{15h^3c^3},
\end{equation}
and the pressure is
\begin{equation}
    P_{\text{rad}}=\frac{1}{3}u=\frac{1}{3}aT^4.
\end{equation}
\end{theory}

\subsection{Adiabatic Expansion of Radiation}
\begin{theory}
For photon gas with $dQ=0$,
\begin{equation}
    T^3V=\text{const}, \qquad PV^{4/3}=\text{const}.
\end{equation}
So the adiabatic index is $\gamma=4/3$.
\end{theory}

\subsection{Planck Spectrum and Stefan--Boltzmann Law}
\begin{theory}
\begin{align}
    u(\nu)d\nu &= \frac{8\pi h\nu^3}{c^3}\frac{1}{e^{h\nu/k_BT}-1}d\nu, \\
    I_\nu &= \frac{2h\nu^3}{c^2}\frac{1}{e^{h\nu/k_BT}-1}, \\
    F &= \sigma T^4,\quad \sigma=\frac{ac}{4}.
\end{align}
\end{theory}

\section{Thermodynamics in Stellar Interiors}

\subsection{Gas and Radiation Pressure}
\begin{theory}
In a fully ionized plasma,
\begin{equation}
    P_{\text{total}} = \frac{\rho k_B T}{\mu m_u} + \frac{1}{3}aT^4.
\end{equation}
Radiation pressure becomes important in massive stars.
\end{theory}

\subsection{Energy Transport Overview}
\begin{intuition}
Stellar heat transport transitions between conduction, radiation, and convection depending on opacity and stability of stratification.
The dominant mechanism is a structural diagnostic for stellar mass and evolutionary stage.
\end{intuition}

\begin{theory}
\begin{itemize}
    \item Mean free path: $\lambda=1/(n\sigma)$.
    \item Conductive flux: $j=-\kappa\,dT/dx$.
    \item Radiative diffusion:
    \begin{equation}
        F_{\text{rad}} = -\frac{4ac}{3\kappa\rho}T^3\frac{dT}{dr}.
    \end{equation}
    \item Convection criterion (Schwarzschild):
    \begin{equation}
        \nabla \equiv \frac{d\ln T}{d\ln P} > \nabla_{\text{ad}}.
    \end{equation}
\end{itemize}
\end{theory}

\begin{extra}
Mass dependence of transport zones (fully convective low-mass stars, solar-type radiative core + convective envelope, massive-star convective cores) can be understood from opacity trends and temperature sensitivity of nuclear burning.
\end{extra}

\section{Eddington Luminosity}
\begin{theory}
Hydrostatic support requires radiative acceleration not to exceed gravity:
\begin{equation}
    L < L_{\text{Edd}}\equiv\frac{4\pi GMc}{\kappa_R}.
\end{equation}
For electron scattering opacity,
\begin{equation}
    L_{\text{Edd}}\approx \SI{1.3e31}{W}\left(\frac{M}{M_\odot}\right).
\end{equation}
\end{theory}
